\chapter{逐篇证据与数值实现档案}

本附录将本专题采用的证据按问题次序列出。旧记录来自已有全文笔记，本轮定向复用不重复计数。新增两篇为电子闭合的观测检验，其中数值部分含混合模型；不能称为新增两次完整双流体模拟。

本专题采用21篇核读证据，其中19篇复用既有记录、2篇新增核读；新增主版本共43个PDF物理页（含封面、参考文献及重复图页），全库144增至146篇。页面数量不等于独立研究数量。


\section{L164：2010年}\label{paper:L164}

\textbf{Slow shocks and conduction fronts from Petschek reconnection of skewed magnetic fields: two-fluid effects}\citep{L164}

\textsf{本轮状态：}已有全文记录复用，不增加计数。

\textsf{核读范围：}PDF 1–45 including all Sections 1–7 and references;  no internal appendix\par\textsf{图式：}Direct PDF visual: Eqs.4–7 (p10),  Eqs.14–18 (p12),  Figs.4–7, 9–14 (pp17, 19, 20, 24, 26, 28, 30, 32, 35, 37). Remaining figures/captions read in text.

\href{run:../../papers/L164.pdf}{本地原文PDF}\quad\href{https://doi.org/10.1088/0004-637X/718/2/1491}{DOI/版本入口}

\paragraph{实际模型与问题}采用准中性、共同密度和共同速度的一维两温碰撞流体描述平行激波，再映射到外部给定的Petschek偏斜磁场重联出流。共同速度来自本问题无电流的特殊约化，并非一般Hall假设。保留离子黏性、两物种经典热导及库仑温度交换；电子惯性忽略，压力梯度电场维持电子力平衡。它不是自洽Hall扩散区模拟。

\paragraph{参数、初边值与实现}例子包含入流马赫数2.4、3.6、5，初始双曲正切速度层宽1–500个归一化离子黏性长度；电子和离子Prandtl数0.013、0.33，交换参数0.019。输运系数随温度的5/2次方变化，交换频率随nTe\textasciicircum{}(-3/2)变化。中心镜面、远端持续入流；稳态用射击积分，时变网格和时间步未报告。偏斜角约≤120°时几何约化较可靠。

\paragraph{求解得到什么}离子黏性子激波先把出流动能转为离子内能，随后长下游区域的碰撞交换加热电子，电子传导又形成前导热波；三个区域及相关时间尺度明显分离。早期离子温度可超过最终平衡值，电子升温滞后，而远下游最终跳跃仍可与单温结果一致。初始速度层足够宽时，局地两温大小关系也可改变，不能压成固定加热比。

\paragraph{成立条件与剩余缺口}原文发现部分强激波早期离子热流达到或超过自由流值，黏性修正甚至违反压力正定条件，故经典闭合恰在关键加热层可能失效。电子热流较小与未包含X点直接电子加热有关，不能外推到所有重联。模型不含自洽Hall场、完整动理学尾、动态磁通管截面和反平行极限，也没有进行实验Ti/Te定量验证。

\paragraph{对Hall比较的意义}这是两温流体可研究热化过程的清楚实例：它把离子流能耗散、物种交换和电子传导在空间与时间上分开，同时主动揭示经典离子闭合的失效。可据此说明“双流体并非不能研究加热”，但不能把给定出流的碰撞激波结果当作无碰撞Hall重联的普适分配律。文中与拾取加热等动理学机制的联系应作为模型边界保留。

\section{L223：2016年}\label{paper:L223}

\textbf{Two-fluid and magnetohydrodynamic modelling of magnetic reconnection in the MAST spherical tokamak and the solar corona}\citep{L223}

\textsf{本轮状态：}已有全文记录复用，不增加计数。

\textsf{核读范围：}PDF 1–16: all body,  methods,  results,  discussion,  figures/captions,  references;  no internal appendix\par\textsf{图式：}Direct PDF visual: Eq.1–5 and definitions (pp3–4),  Fig.3 (p7),  relaxation equations/Figs.4–8 (pp9–10, 13–14);  remaining figure captions read in full text

\href{run:../../papers/L223.pdf}{本地原文PDF}\quad\href{https://doi.org/10.1088/0741-3335/58/1/014041}{DOI/版本入口}

\paragraph{实际模型与问题}采用二维直角和紧环径比轴对称环面Hall MHD描述同向电流磁通绳合并，保留电子漂移、电子标量压力、电阻、超电阻及热传导。前列方程仍演化总压，图3才给独立离子、电子温度的初步扩展：离子由黏性供热，电子由电阻和超电阻供热。后半篇的Taylor弛豫计算只估计磁能释放，不能混作物种加热模拟。

\paragraph{参数、初边值与实现}基准密度5×10\textsuperscript{18} m\textsuperscript{-3}、长度1 m、磁场0.5 T，初始Ti=Te=10 eV，归一化离子惯性长度0.145。典型电阻10\textsuperscript{-5}、离子黏性10\textsuperscript{-3}、平行电子热导0.1、垂直离子热导10\textsuperscript{-7}；图3超电阻10\textsuperscript{-8}。使用HiFi及前序模型，但本文没有报告该两温加热图专用的网格、时间步或收敛检查，完整两温方程亦未列出。

\paragraph{求解得到什么}图3的初步两温环面计算产生下游离子双热峰，峰值可达keV量级，与早期MAST高驱动观测有定性联系；电子热源主要跟随超电阻涉及的电流梯度，分布不同于离子。中心或空心电子温度与岛排出的联系仍是初步解释。弛豫模型另得到几何相关磁能释放及近似极向磁场平方标度，但不能给出热量物种比例或空间分布。

\paragraph{成立条件与剩余缺口}超电阻是自由模型参数，并直接影响电子加热位置；有限标量闭合不描述完整粒子分布或非热尾。离子热峰高于后来较低驱动实验，实验比较仍属初步，不能称为加热比例定量验证。Stanier2014学位论文、源代码和独立数据未核读，本文也未给两温热图的专门收敛证据；没有推导普适Qi/Qe。

\paragraph{对Hall比较的意义}本篇直接证明两温Hall扩展能够在指定耗散、传导和几何条件下计算离子与电子局地热图。证据应表述为“模型能做条件性热预算和实验比较”，不能表述为“Hall项本身提供不可逆热化”或“物种分配已经普遍解决”。同时必须区分前列总压方程、后续初步两温计算和独立磁能弛豫估计三种证据。

\section{Kaptanoglu2020：2020年}\label{paper:Kaptanoglu2020}

\textbf{Two-temperature effects in Hall-MHD simulations of the HIT-SI experiment}\citep{Kaptanoglu2020}

\textsf{本轮状态：}已有全文记录复用，不增加计数。

\textsf{核读范围：}PDF 1–35: all body,  methods,  results,  conclusions,  references and appended repeated/enlarged figure pages;  no internal appendix.\par\textsf{图式：}Direct PDF visual: Eq.1–13/Table I (pp3–4),  Figs.2–4 (pp7–9),  Figs.5–13/scan results (pp11–17);  overview of appended repeated images pp19–35. Key Ti/Te map and heating power plot inspected. Some noncentral scan figures reviewed at contact-sheet scale.

\href{run:../../papers/Kaptanoglu2020_accepted.pdf}{本地原文PDF}\quad\href{https://doi.org/10.1063/5.0006311}{DOI/版本入口}

\paragraph{实际模型与问题}在PSI-Tet和NIMROD中采用三维两温Hall MHD，独立演化Ti、Te，欧姆定律保留Hall漂移及增强电子质量的简化电流时间导数惯性项。电子温度随电子流对流和压缩；电子得欧姆热、离子得黏性热，两者通过库仑交换及各向异性热传导联系。闭合仍是标量压力、Spitzer电阻和常数各向同性黏性，另有人工粒子扩散。

\paragraph{参数、初边值与实现}注入频率14.5–68.5 kHz，边界密度0.75×10\textsuperscript{19} m\textsuperscript{-3}、温度3 eV，黏性550 m\textsuperscript{2}/s、扩散250 m\textsuperscript{2}/s，有效质量比100。PSI-Tet用103887四面体、2.8 cm三阶单元，等效9 mm，Crank–Nicolson步长≤40 ns；NIMROD用1.8 cm四阶极向单元和11环向模。两代码注入器几何不同，另扫描壁温、密度、扩散和增强质量。

\paragraph{求解得到什么}独立两温改变了单温模型的温度、密度与电流演化；较高驱动频率下离子明显热于电子。图4显示两物种相似的大尺度空间结构，但温度幅度不同；图3列各物种热源、交换和壁面热流，积分范围包括注入器。作者称与离子Doppler和初步Thomson数据定性一致，但明确未做严格定量验证。压缩功也不能全部视为不可逆热量。

\paragraph{成立条件与剩余缺口}温度依赖壁面条件和人工扩散，尚未达到对扩散系数不敏感的收敛极限；真实驱动电路、中性/原子/辐射过程及更真实黏性仍缺失。已有勘误DOI10.1063/5.0032632经元数据确认，但更正全文未取得；本次只保证35页接收稿完整核读，不宣称掌握更正后的全部公式或精确比值。两个代码的差异也不全来自两温处理。

\paragraph{对Hall比较的意义}这是本分支取得的最直接装置尺度两温Hall实例，证明可同时计算独立温度、局地热图和带输运损失的物种预算。但结果受驱动、边界与耗散闭合约束，且部分功率来自注入器反转过程，并非仅重联X点。应采用定性和条件性结论；未核勘误、人工扩散敏感性及缺失电路模型使其不能支持参数无关的Qi/Qe或精密实验验证。

\section{L247：2019年}\label{paper:L247}

\textbf{High-order two-fluid plasma solver for direct numerical simulations of plasma flows with full transport phenomena}\citep{L247}

\textsf{本轮状态：}已有全文记录复用，不增加计数。

\textsf{核读范围：}PDF1–44: entire text and references including Appendix A plasma transport,  B quasineutrality,  C single-fluid limits.\par\textsf{图式：}Direct visual Eq15–19 (p11),  identical-temperature restriction (p16),  Figs1, 3, 4, 7, 9, 10, 11 (pp18, 21, 25, 27, 35),  AppendixA EqsA1–A9, A15–A19 (pp31, 33). All other figures/captions read in text.

\href{run:../../papers/L247.pdf}{本地原文PDF}\quad\href{https://doi.org/10.1063/1.5082190}{DOI/版本入口}

\paragraph{实际模型与问题}构建准中性、忽略电子惯性与位移电流的Braginskii两流体方程，保留独立物种内能、磁场方向相关黏性和热流、摩擦、热力、电阻加热及温度交换；附录列出输运系数和单流体极限。但第IV节明确四项实际演示均强制Ti=Te，并删去独立离子能量演化，实际测试的是单温扩展Hall系统，不能当作两温分配验证。

\paragraph{参数、初边值与实现}波测试用384点；激波Rei为6.3×10\textsuperscript{3}和2.3×10\textsuperscript{4}。重联di/L=1、Rei约124–850、Rem约58–890，主例426/112；参考长度10 m、密度10\textsuperscript{15} m\textsuperscript{-3}。CFDNS用六阶紧致差分和物理输运；主例约256×512收敛，以768×1536为参考，观测阶数6.58。正文未给具体时间步，另有更高Re和激波网格扫描。

\paragraph{求解得到什么}在所选物理输运系数下，波色散、激波和重联基准得到高阶网格收敛；局地共同温度图显示磁能转化及热扩散的重要性。激波误差阶数约6.08–6.14，重联磁通在所试较大黏性和磁雷诺数处趋于弱敏感。论文提供了完整物种热方程及物理正则化的实现路径，但没有给出独立Ti/Te、交换过程或加热比例的验证结果。

\paragraph{成立条件与剩余缺口}最重要边界是所有验证算例固定Ti=Te，不能从完整方程形式推断已演示双温预测。经典碰撞输运在冠层扩散层可能不满足局地平衡排序；磁通及饱和时间随网格或Re收敛，也不能证明微观热化和物种分配正确。零电子惯性等约化排除相关电子尺度；原始输运推导、CFDNS源代码及独立双温测试未复核，也无实验温度比对照。

\paragraph{对Hall比较的意义}应将本篇用作能量方程、各向异性输运和数值收敛的方法证据，尤其支持区分物理输运与隐含数值耗散。作者关于高Re宏观结果可代表无碰撞重联的讨论，只由其实际磁通等指标检验，不能扩张为无碰撞加热闭合获验证。全文44页含附录均已读；正式出版年为2019，2018是预印本日期，两者不是独立文献。

\section{Wang2015：2015年}\label{paper:Wang2015}

\textbf{Comparison of multi-fluid moment models with particle-in-cell simulations of collisionless magnetic reconnection}\citep{Wang2015}

\textsf{本轮状态：}已有全文记录复用，不增加计数。

\textsf{核读范围：}1-15（arXiv v3全部）\par\textsf{图式：}图1-10及式20-24、36；图4、8-10视觉核验

\href{run:../../papers/Wang2015.pdf}{本地原文PDF}\quad\href{https://doi.org/10.1063/1.4906063}{DOI/版本入口}

\paragraph{实际模型与问题}五矩和十矩全Maxwell，局地Hammett-Perkins松弛；k0e=k0i=1/de为targeted，10\textasciicircum{}4/de近五矩。

\paragraph{参数、初边值与实现}主Harris例Lx×Ly=100×50di, a=.913di, nb=.3, Ti/Te=5, mi/me=25, vA/c=1/15, 扰动.1；流体4095×2047（约8格/de），PIC5760×2880（约12格/de），每物种2.7e9粒子。另有mi/me=25, 100, 256, 400质量比扫描。

\paragraph{求解得到什么}图4通量时间史明显不同；同重联磁通下targeted十矩的X点附近电子流和非对角压强较接近PIC，五矩及强松弛出现PIC没有的磁岛链。质量扫描给出电子层长约(me/mi)\textasciicircum{}(3/8)、厚约(me/mi)\textasciicircum{}(1/3)。

\paragraph{成立条件与剩余缺口}对比结构常取不同时间（图10 PIC75、流体132.5）；Pxy峰值PIC约.024对流体.006，标量化仍偏强；下游欧姆项峰位不一致、抵消不足。真实质量比层长12-17de和厚.4-.6de是外推，不是直接计算；无完整热流张量检验。

\paragraph{对Hall比较的意义}能复现部分局地压力支撑，不能写成重联全过程或所有压力分量定量吻合；局地闭合参数仍需选择。

\section{Ng2015：2015年}\label{paper:Ng2015}

\textbf{The island coalescence problem: Scaling of reconnection in extended fluid models including higher-order moments}\citep{Ng2015}

\textsf{本轮状态：}已有全文记录复用，不增加计数。

\textsf{核读范围：}1-8（全部）\par\textsf{图式：}图1-6文字及图注；主要结果图已视觉复核

\href{run:../../papers/Ng2015.pdf}{本地原文PDF}\quad\href{https://doi.org/10.1063/1.4935302}{DOI/版本入口}

\paragraph{实际模型与问题}电阻MHD、Hall MHD电阻／超电阻、十矩局地松弛对照PIC；Fadeev岛合并。

\paragraph{参数、初边值与实现}epsilon=.4, nb=.2n0, Lx×Lz=4piλ×2piλ, x周期, z导体／反射；mi/me=25，Ti=Te，beta=1；平均窗0-1.5tA，模拟至2.5tA。Hall eta=.009或etaH=1e-4及eta=1e-5；mu=.01。

\paragraph{求解得到什么}Hall区峰值／平均率对系统尺寸仅弱敏感，约λ\textasciicircum{}-0.2至\textasciicircum{}-0.25；十矩调低ki到1/(3di)改善小系统离子压力支撑，但大系统次级岛不断抛出，合并完成且缺少PIC弹回。

\paragraph{成立条件与剩余缺口}本文比较PIC平均率约λ\textasciicircum{}-0.8；引言早期研究的\textasciicircum{}-0.5不可混用，盒子宽高比不同。十矩ki是可调参数，出现PIC中没有的次级岛；本文不是十矩全面成功的证据。

\paragraph{对Hall比较的意义}必须比较同几何、同平均窗；Hall率可高2-3倍。

\section{Ng2017：2017年}\label{paper:Ng2017}

\textbf{Simulations of anti-parallel reconnection using a nonlocal heat flux closure}\citep{Ng2017}

\textsf{本轮状态：}已有全文记录复用，不增加计数。

\textsf{核读范围：}1-5（全部）\par\textsf{图式：}图1-5文字及图注；主要结果图已视觉复核

\href{run:../../papers/Ng2017.pdf}{本地原文PDF}\quad\href{https://doi.org/10.1063/1.4993195}{DOI/版本入口}

\paragraph{实际模型与问题}全Maxwell十矩；非磁化三维Hammett-Perkins热流Fourier闭合，chi=sqrt(4/(9pi))；有限体积输运、局地隐式源与FFTW。

\paragraph{参数、初边值与实现}Fadeev岛、epsilon=.4, nb=.2, Lx×Lz=4piλ×2piλ, mi/me=25, vt, e/c=.35, Ti=Te, beta=1, ωpe/Ωce=2，10\%磁扰动；平均0-1.5tA。

\paragraph{求解得到什么}非局地闭合得到1.20di片宽，PIC1.15di，Hall .18di；离子压力支撑区与速度、非旋转对称度改善；平均率λ\textasciicircum{}-0.45比局地\textasciicircum{}-0.2更接近PIC\textasciicircum{}-0.8。

\paragraph{成立条件与剩余缺口}仍未复现PIC尺寸标度，离子非旋转对称度偏低；次级岛影响速率；电子各向异性偏弱；非磁化闭合不应直接外推强导引场。Langmuir阻尼测得.044ωpe，精确Vlasov .034ωpe，存在约29\%误差。

\paragraph{对Hall比较的意义}成功是特定离子结构改善与减少调参，不是所有动力学完全重现。

\section{L256：2018年}\label{paper:L256}

\textbf{Temperature gradient driven heat flux closure in fluid simulations of collisionless reconnection}\citep{L256}

\textsf{本轮状态：}已有全文记录复用，不增加计数。

\textsf{核读范围：}期刊PDF1-16全部\par\textsf{图式：}图1-10及表1；图1、3-6、8-10与表1视觉复核

\href{run:../../papers/L256.pdf}{本地原文PDF}\quad\href{https://doi.org/10.1017/s002237781800048x}{DOI/版本入口}

\paragraph{实际模型与问题}全Maxwell十矩、对压力偏离各向同性的部分作Laplacian闭合，不是直接扩散完整温度；以Vlasov热流高阶矩作对照。

\paragraph{参数、初边值与实现}GEM质量比25、256×128基准及1024×512加密；另流体质量比400/1836。k0s分别为GEM3/ds、100×50di WHBG的1/(3ds)、岛合并1/(2ds)；不是一组参数全场景。岛lambda5/10/15/25di，两套网格512×256到2048×1024。

\paragraph{求解得到什么}原松弛闭合的热流散度形状和符号在某些区域错误；梯度闭合改善上游假热流与电子非对角压力。岛合并平均0-1.5tA率标度约lambda\textasciicircum{}(-.73)，出现反弹且无二次岛。表1lambda5至25平均率.180→.052；两套网格差小。GEM高分辨梯度流体后期斜率近Vlasov，但启动仍迟；高分辨局地松弛反而更慢。

\paragraph{成立条件与剩余缺口}k0按与动理学一致调节，尺寸标度也依赖k0；不能宣称无参预测。实际质量比只展示流体结构，不是同质量比完整动理学验证。显式Laplacian子循环中不交换块边界、不更新速度密度，作者承认块边界误差。GEM温度条件缺常规1/2，岛平衡坐标与边界文字不一致，复现需输入核对。

\paragraph{对Hall比较的意义}重要改进是直接检查热流矩再验证非线性反弹；其成功同时有问题相关的校准与数值限制。

\section{Ng2020：2020年}\label{paper:Ng2020}

\textbf{An improved ten-moment closure for reconnection and instabilities}\citep{Ng2020}

\textsf{本轮状态：}已有全文记录复用，不增加计数。

\textsf{核读范围：}1-8（作者稿全部，含附录A）\par\textsf{图式：}图1-7、式4-6、附录A；2026-11-14定向视觉复核全部结果图及附录页

\href{run:../../papers/Ng2020.pdf}{本地原文PDF}\quad\href{https://doi.org/10.1063/5.0012067}{DOI/版本入口}

\paragraph{实际模型与问题}十矩全Maxwell温度梯度热流，与局地各向同性化和非局地无磁化闭合比较；FV输运、局地隐式源、超时间步抛物子步。

\paragraph{参数、初边值与实现}Fadeev与Ng2015同系；lambda5di用640×320，同dx扫描系统尺寸，mi/me25, k0s=1/ds；导引场0-.4-1-2等。LHDI另例mi/me36, Ti/Te10, L=rhoi, vt, e=.06c, 盒6.4L×12.8L, 模数8。

\paragraph{求解得到什么}平均率随岛尺寸下降可接近PIC约lambda\textasciicircum{}-0.8；岛合并最大率随Bg升到1先升后降，不能套用Harris片单调下降。离子宽层及非旋转对称度改善。相同k0=1/ds可同时激发重联与LHDI。

\paragraph{成立条件与剩余缺口}LHDI增长率默认1.33Ωci、调k为ky后.86，动力学1.1，局地闭合.34，非局地.84；仍非完全吻合。电子各向异性空间结构仍不准、跨场热流过大；未研究LHDI对重联的非线性作用。附录指出零热流、无碰撞的对称稳态停滞点会驱动不物理Pxx趋零。 图6的流体电场图为显示除以2，不能仅由色彩图形近似声称幅度吻合。

\paragraph{对Hall比较的意义}梯度闭合并非更高阶就自动正确；改善来自k\textsuperscript{2}/|k0|的尺度选择，仍有自由参数和磁化缺失。

\section{TenBarge2019MMS：2019年}\label{paper:TenBarge2019MMS}

\textbf{An Extended MHD Study of the 16 October 2015 MMS Diffusion Region Crossing}\citep{TenBarge2019MMS}

\textsf{本轮状态：}已有全文记录复用，不增加计数。

\textsf{核读范围：}1-17（全部）\par\textsf{图式：}PDF5、6、9、10、11页直接视觉核验式12、图1/6-10；17页全文及13图图注读完。正式HTML新增第14图未视觉核验。

\href{run:../../papers/TenBarge2019MMS.pdf}{本地原文PDF}\quad\href{https://doi.org/10.1029/2019JA026731}{DOI/版本入口}

\paragraph{实际模型与问题}两物种全Maxwell五矩/十矩Gkeyll，局地Hammett–Perkins型压力松弛；直接有限矩流体证据。

\paragraph{参数、初边值与实现}mi/me100, c/vAs500；2D40.96×20.48di, 2048×1024；3D40.96×20.48×10.24di, 1024×512×256；周期双片；Bg/Bsx.099, Bmx/Bsx-1.7, Tem/Tes12.3, Tim/Tis5.63；k0s ds=.1，种子.1，磁噪声.0005。

\paragraph{求解得到什么}图1直接给离子/电子热能和流能，十矩2D/3D全域预算相近、离子热增量较大；电子各向异性峰约2.5近MMS幅度但位置在磁鞘侧而观测在磁层侧；局地闭合抑制LHDI，剪切/kink仍出现。

\paragraph{成立条件与剩余缺口}热能增量不是物种熵产Qi/Qe；五矩过快LHDI并明显总能损失，不可信作分配预测；未给网格/质量比独立收敛；nm/ns=16.7与低密度磁层文字冲突，l0=2π/k0≈d与k0d=.1冲突；不静默修正输入。

\paragraph{对Hall比较的意义}模型身份、温度/能量/熵产区别及跨文献判断见独立加热专题；不把本篇单独作为普适分配定律。

\section{E552：2021年}\label{paper:E552}

\textbf{Fluid simulations of three-dimensional reconnection that capture the lower-hybrid drift instability}\citep{E552}

\textsf{本轮状态：}已有全文记录复用，不增加计数。

\textsf{核读范围：}arXiv2008.06440v3 1-20全部，对应2021年期刊\par\textsf{图式：}图1-9；闭合式、图4、6-9视觉复核

\href{run:../../papers/E552.pdf}{本地原文PDF}\quad\href{https://doi.org/10.1017/s0022377820001683}{DOI/版本入口}

\paragraph{实际模型与问题}全Maxwell五矩与十矩比较；压力偏离各向同性的梯度闭合chi3, k0s1/ds；Vlasov用于二维电流片不稳定性对照。

\paragraph{参数、初边值与实现}CWENO+RK3，Maxwell FDTD，Laplacian二阶中心，每子循环交换压力边界，三维用8次；压力对角下限。三维100×25×6di、2400×600×150、mi/me100, c/VA30，Harris a.5, nb.2, Ti/Te5；磁通种子.1、电子密度噪1e-8。独立静电LHDI nb.001、Ti/Te10、a约.67、质量比36/100/250至流体1836。

\paragraph{求解得到什么}静电LHDI质量比100时Vlasov/十矩梯度/五矩/十矩松弛增长率2.2/2.4/3.2/.67 Omega\_i；最快模随质量比变化，固定单一模比较会误导。电磁支Vlasov约.41、流体.21，均出现超过B0的磁扰动及强kink。三维较二维启动早、后期片变宽，出流湍动强而电子层相对适中；密度噪1e-10至1e-2使LHDI启动约30→4 Omega\_i\textasciicircum{}-1。

\paragraph{成立条件与剩余缺口}高质量比无同参数Vlasov验证；电磁支结构、启动与增长仍不同，异常高饱和原因未解释。CSSI只是形态相似的候选，未确诊。三维率是z平均电场代理；电子de约2.4格，未给全三维网格收敛。对角下限不等于完整压力正定保证；印刷Lap约+k0\textsuperscript{2}应理解为负号Fourier符号。

\paragraph{对Hall比较的意义}该工作同时检验重联与不稳定性并强调噪声/启动竞争，优于仅比较峰值；但不能把3D形态相近称为完整动理学等价。

\section{Closure2026：2026年}\label{paper:Closure2026}

\textbf{Capturing Secondary Kinetic Instabilities in Three‐Dimensional Dayside Reconnection Using an Improved Gradient‐Based Closure}\citep{Closure2026}

\textsf{本轮状态：}已有全文记录复用，不增加计数。

\textsf{核读范围：}1-23（arXiv2602.16960v1全部，含附录A）\par\textsf{图式：}图1-11、附录A1及式A1-A12；2026-11-14定向视觉复核全部结果图及附录页

\href{run:../../papers/Closure2026.pdf}{本地原文PDF}\quad\href{https://doi.org/10.1029/2026ja035250}{DOI/版本入口}

\paragraph{实际模型与问题}Gkeyll三维非对称十矩；温度梯度闭合与局地闭合对照；顶点对称热扩散和alpha=.75限制器。

\paragraph{参数、初边值与实现}盒40.96×20.48×10.24di，全周期双电流片，1152×576×288，mi/me100, c/vA200, Bg/Bsx=.099；nm/ns=.062, Bxm/Bxs=-1.696, Tem/Tes12.3, Tim/Tis5.63；k0αdα约1，磁噪声.0005，psi种子.1。

\paragraph{求解得到什么}梯度闭合出现约5Ωci\textasciicircum{}-1内LHDI、约20后kink与通量绳/次级岛，局地同参不出现；三维率先峰后降，压力张量支撑X点非理想场。梯度闭合成本约3倍。 图6右（PDF11页）中，梯度闭合相对局地闭合的离子压力能及流能变化更小，电子压力能增长较快；这不等于已测得更大不可逆电子熵产。

\paragraph{成立条件与剩余缺口}温度各向异性峰值侧别与MMS/PIC不符，非旋转对称度与各向异性区域过宽，作者认为背景跨场扩散产生不物理各向异性。独立2D Vlasov对照kink仅数个离子周期出现，流体需数十周期，且遗漏小尺度模。正文未给质量比/网格独立收敛序列。

\paragraph{对Hall比较的意义}证明某些先前受抑制的横向不稳定性可在十矩中出现，不是准确恢复全部三维动理学。建议的磁场方向热流及弱各向同性化仍为未来改进。

\section{L244：2018年}\label{paper:L244}

\textbf{A reduced Landau-gyrofluid model for magnetic reconnection driven by electron inertia}\citep{L244}

\textsf{本轮状态：}已有全文记录复用，不增加计数。

\textsf{核读范围：}HAL02049658v1封面及正文1–26全读；物理页1–27\par\textsf{图式：}式2.1–4.21、Hamiltonian/Casimir推导及图1–7；阻尼算子、图3–6视觉复核

\href{run:../../papers/L244.pdf}{本地原文PDF}\quad\href{https://doi.org/10.1017/S002237781800051X}{DOI/版本入口}

\paragraph{实际模型与问题}从Scott型gyrofluid按epsilon低频、delta\textsuperscript{2}=me/mi、beta\_e\textasciitilde{}delta\textsuperscript{2}约化为六场；含离子FLR、电子双温、电子惯性和非局地Landau闭合。实际重联只算冷离子、Ni=Tperpi=0、Tperpe恒零的子模型。

\paragraph{参数、初边值与实现}delta=.1, beta\_e=.05, lambda3；线性半盒Lx46, Ly24pi, Lz96pi；非线性ky1/12, kz1/48，Ly12pi, Lz48pi。显式AB3、y/z并行；非线性热流刚性使步长至多缩小100倍，改用单螺度坐标y*=y+(kz/ky)z的等价二维推进。全文未列具体格点数及绝对步长。

\paragraph{求解得到什么}Landau算子用局地平行梯度和磁场方向二阶空间平均的Fourier分母，非简单固定k0扩散。阻尼降低平行温度扰动、线性增长率，高m斜模可稳定；同一gamma*t上岛面积曲线相似，意味着真实时间中发展变慢，而非未受影响。饱和岛大小近似不变，细密度/电流/温度结构变平滑。闭合去除自由能，未出现动理学参考中的重联—加热延迟。

\paragraph{成立条件与剩余缺口}斜单螺度不是多螺度三维湍流；没有热离子完整六场非线性验证。无Landau时数值能损约总能量.2\%，却约为同期磁能释放20\%，不可忽略。未出现延迟可能是闭合过度耗散，作者提出高阶闭合问题。自由能损失解释为未解析相混合/电子加热，非直接分布函数测量。

\paragraph{对Hall比较的意义}给出可推导、可检验的闭合成功与失败并存例子；必须同时检验时序、饱和、能量预算，不能以一项吻合宣称动理学等价。

\section{F635：2026年}\label{paper:F635}

\textbf{Counter-streaming heat-flux closure for electron-only collisionless magnetic reconnection}\citep{F635}

\textsf{本轮状态：}已有全文记录复用，不增加计数。

\textsf{核读范围：}arXiv2607.07106v1全部10页；公开三个输入文件已下载核对关键设置\par\textsf{图式：}图1-6全部视觉核查；特别核对图4最佳时间指标、图5独立色标

\href{run:../../papers/F635.pdf}{本地原文PDF}\quad\href{https://doi.org/10.48550/arxiv.2607.07106}{DOI/版本入口}

\paragraph{实际模型与问题}2026预印本：固定离子背景，两个电子十矩流体；对照单电子十矩及OSIRIS PIC。不是通常电子—离子双流体。

\paragraph{参数、初边值与实现}48de×64de双周期，480×640；力自由双层半厚de，Brec=.5、Bg/Brec=.1、beta\_rec=.1。PIC每格625粒子、二阶形函数、dt=.066667/omega\_pe；流体Gkeyll分维有限体积CFL2/3、k=.1/de各向同性松弛。两电子群按片上下来源标记，tanh分配宽de、密度下限.03n0；初始速度相同。Zenodo输入中plasmaBetae0=.3等旧声明不参与实际温度初始化；实际vth=.111803与Brec=.5给beta=.1。

\paragraph{求解得到什么}两群相对漂移与不同压力可重建总三阶中心矩的大部分结构，改善单流体的密度空洞与局部压力堆积；热流符号/形状更接近PIC。在示例中单群内部热流对局部Qxxx贡献约4\%；但双流体重建热流幅度仍偏高。

\paragraph{成立条件与剩余缺口}图4是每个PIC参数算例在时间内取最大相似度，非双流体全参数性能或时间平均；图5流体色标不同，不能直接认定幅度一致。高beta时人为分群梯度引发自由膨胀，晚期同群再度多束化；三维、动态分裂和离子耦合均未验证。156倍是按粒子相空间变量数与20矩的表示压缩估算，不是实测运行加速。仍有弱松弛而非严格Q=0。

\paragraph{对Hall比较的意义}有价值的新方向是利用来源分群减少每群的闭合负担；不能推成单电子流体在数学上不可能表达这些矩，或热流闭合已解决。

\section{Wetherton2019：2019年}\label{paper:Wetherton2019}

\textbf{Validation of Anisotropic Electron Fluid Closure Through In Situ Spacecraft Observations of Magnetic Reconnection}\citep{Wetherton2019}

\textsf{本轮状态：}新增全文核读。

\textsf{核读范围：}OSTI接收稿PDF1–20全部；封面、正文、方法、讨论、参考文献、后部重复图页均读完，无文内附录。\par\textsf{图式：}图1–3、式1–6；直接查看PDF4、6、8、13、15、17–20页，图2密度转折与图3截面特别核对。

\href{run:../../papers/Wetherton2019.pdf}{本地原文PDF}\quad\href{https://doi.org/10.1029/2019GL083119}{DOI/版本入口}

\paragraph{实际模型与问题}各向异性电子EoS的MMS观测检验；H3D数值部分为动理学离子加流体电子，非完整两流体。外部来源、俘获/通过粒子决定p\_parallel(n, B)与p\_perp(n, B)，未演化全非旋转对称压力。

\paragraph{参数、初边值与实现}2015-10-31强导场近对称MMS3事件，轨迹距X点约7–10di。来源n=7.1cm\textasciicircum{}-3、B=55.5nT、Tperp=32.4eV、Tparallel/Tperp=1.22。H3D力自由初态，Bg/Br=1.75、总场beta\_e=.02、Ti/Te=9.1；80×40di、2048×1024网格、400离子粒子/格；显式电阻破冻结。t=80/Omega\_ci选宽约3.5di截面对照；本文未列完整边界、时间步、电阻值或收敛序列，需前序方法与输入。

\paragraph{求解得到什么}观测平行温度在低密度近Boltzmann、高密度近CGL，EoS跨越转折；混合模拟避免低质量比100 PIC在密度空洞处的过强平行降温。强导场力平衡加绝热离子给归一化电子温升弱beta依赖，典型条件约.02；这是温升而非两物种不可逆热量比。

\paragraph{成立条件与剩余缺口}数值验证属混合，不能计为完整双流体。来源与截面匹配使其不等同无参数预测；未给统一残差/仪器协方差。峰值B≈Bg、Ti\ensuremath{\gg}Te及绝热离子限制标度；图2c横轴电子重联场beta，式4–6中beta\_r为离子beta，不能混用。式1–2周围平行热流符号与三阶矩定义含二分之一规范疑点，报告采用q\_a=Q\_abb/2。前序EoS推导、原始MMS数据和方法论文未完整重读。

\paragraph{对Hall比较的意义}强于单峰温度符合的证据是跨密度区间的各向异性转折及分布边界。支持电子闭合，不能验证离子热化或通用分配。相对标量Hall新增的是压力自由度与受俘获限制的热流响应；可以嵌入各向异性Hall，非对所有Hall的否定。

\section{Wetherton2021：2021年}\label{paper:Wetherton2021}

\textbf{Anisotropic Electron Fluid Closure Validated by in Situ Spacecraft Observations in the far Exhaust of Guide-field Reconnection}\citep{Wetherton2021}

\textsf{本轮状态：}新增全文核读。

\textsf{核读范围：}OSTI接收稿PDF1–23全部；正文、模型、全部图表、讨论、参考文献读完，无文内附录。\par\textsf{图式：}图1–10与表1、式1–4；直接看PDF6–7、9、11–12、14–19页，重点图6–7的偏差区和图9–10的反演增能。

\href{run:../../papers/Wetherton2021.pdf}{本地原文PDF}\quad\href{https://doi.org/10.1029/2020JA028604}{DOI/版本入口}

\paragraph{实际模型与问题}MMS远下游观测与电子各向异性EoS比较；数值取矩求压力而非完整双流体重联计算。论文第2节复用PIC解释密度与电势，不能当作新增纯流体模拟。

\paragraph{参数、初边值与实现}主事件2016-01-21约01:06UT，MMS2，距X线约100di、Bg/Br约.7、略非对称；以2015-10-31近区事件为对照。分别选择入流1/出流/入流2来源n=47/41/40cm\textasciicircum{}-3，B=30/29/41nT，Tperp=28/30/30eV，Tparallel/Tperp=.9/1/.9以匹配温度。各区域数百采样，但四个来源参数不是独立直接远端测量。复用PIC图1质量比1836、20×20di；本文不提供新流体网格收敛。

\paragraph{求解得到什么}三来源EoS大体解释平行与垂直温度随局地n、B的变化，给出三条不同区域特征曲线。部分垂直温度明显超过预测，作者联系到区域磁压缩、局地俘获和有限通过时间，以DeltaE=mu积分partial\_t B dt反演额外增能；深俘获例1.35与B/Binf≈1.34相符，出流额外项通常小于.1。

\paragraph{成立条件与剩余缺口}来源参数按观测选取；反演DeltaE亦使用观测，不能视作独立预测。原文明确DeltaE没有对n、B的形式函数，未形成通用新EoS。压缩起源、是否延伸至X线不明；局地磁场极小也不必总有大增能。未给通用拟合残差、置信区间或完整诊断响应。图7b曲线有0/10/20eV人为偏移，不可读为绝对温差。前序时序定位和仪器论文、原始MMS数据未复算。

\paragraph{对Hall比较的意义}把近X线闭合支持扩展到远区同时给出失败边界。不能把经三来源拟合的吻合当成任意条件下预言；额外垂直增温需要历史或新增状态信息。对标量Hall是各向异性约束，对既有EoS-Hall则是局地闭合与来源假设的限制；未解决离子电子不可逆热化比。

\section{Tanabe2015MAST：2015年}\label{paper:Tanabe2015MAST}

\textbf{Electron and Ion Heating Characteristics during Magnetic Reconnection in the MAST Spherical Tokamak}\citep{Tanabe2015MAST}

\textsf{本轮状态：}已有全文记录复用，不增加计数。

\textsf{核读范围：}所记录版本全部5个PDF物理页；正文、方法、结果、讨论、图表、参考文献及适用文内附录均核读。\par\textsf{图式：}直接查看PDF物理页2, 3, 4；All Figs1–5;  experimentgeometry, Te/ne/Titomography, timeevolution, guidefieldscan

\href{run:../../papers/Tanabe2015MAST.pdf}{本地原文PDF}\quad\href{https://doi.org/10.1103/PhysRevLett.115.215004}{DOI/版本入口}

\paragraph{实际模型与问题}本文是球形托卡马克的强导场双环合并实验，以汤姆孙散射测量电子温度和密度，以离子多普勒层析重建温度空间分布，并与已有双流体模拟比较。实验同时讨论电子与离子的碰撞能量交换、各向异性热输运及封闭磁面约束。论文自身不进行独立霍尔模拟，温度结构的符合是对具体模型解释的实验约束。

\paragraph{参数、初边值与实现}环向导场约为 0.3 至 0.8 特斯拉，重联场约为 0.07 至 0.15 特斯拉；离子皮肤深度约 0.1 米，离子回旋半径小于 0.01 米。电子诊断分别为 300 和 130 通道，空间分辨率约 15 与 10 毫米；离子诊断为 32 弦碳离子光谱层析。二维剖面组合多个放电脉冲和测量位置，重联场依据合并后磁平衡重建估计。

\paragraph{求解得到什么}电子在重联中心附近形成宽约 2 至 5 厘米的热点，温度在小于 1 毫秒内由约 10 升到 200 电子伏；离子温升分布在出流及下游停滞区。经过约 4 至 10 毫秒的两物种能量交换时间，两个温度剖面均形成三峰。更高导场使中心电子热点更尖，但所测下游离子温度基本不变，与较弱导场机制的趋势不同。

\paragraph{成立条件与剩余缺口}原文明确指出缺乏直接磁探针测量限制了热点成因的判别。黏性、压缩阻尼与输运解释依赖模型，温度图并不是完整局地功率或熵产预算；杂质离子温度与主离子温度的对应也依赖诊断和碰撞条件。全部二维点并非单次放电同时测得；独立诊断论文、学位论文和原始脉冲数据未在本次重新分析。

\paragraph{对Hall比较的意义}这是双温流体建模线的重要实验支撑：原文直接引用 L223，并将下游温升与含碰撞黏性耗散的双流体结果联系。但热点位置相似不足以唯一验证闭合，强导场下约束增强和驻留时间延长也会改变最终温度。该文支持具体装置中的局部加热研究，没有给出普适物种分配比；正式版 5 页及全部五幅图已核读。

\section{F674：2015年}\label{paper:F674}

\textbf{Study of energy conversion and partitioning in the magnetic reconnection layer of a laboratory plasma}\citep{F674}

\textsf{本轮状态：}已有全文记录复用，不增加计数。

\textsf{核读范围：}所记录版本全部11个PDF物理页；正文、方法、结果、讨论、图表、参考文献及适用文内附录均核读。\par\textsf{图式：}直接查看PDF物理页6, 7, 8, 9；Figs6–8,  TableI,  Eqs(2)–(6)

\href{run:../../papers/F674_published.pdf}{本地原文PDF}\quad\href{https://doi.org/10.1063/1.4920960}{DOI/版本入口}

\paragraph{实际模型与问题}本文通过磁重联实验装置的局部测量建立有限控制体能量收支，并与装置几何下及开放边界下的全动理学粒子模拟比较。分别积分电子和离子的电场功，再区分有序流能、热储能与边界输运。它是实验能量预算和动理学对照研究，不能列为霍尔流体独立预测温度比例的数值证据。

\paragraph{参数、初边值与实现}本次读取 2015 年正式版 11 页，保留的 2014 年预稿按同一论文处理。开放粒子模拟采用真实质子质量比 1836，初始离子与电子温度比为 5，背景密度比为 0.23，背景电子温度比为 0.76。模拟盒为 20 乘 20 个参考离子皮肤深度，网格为 5120 乘 5120，每种粒子约一百亿个；电流片半宽为 0.5 个参考离子皮肤深度。

\paragraph{求解得到什么}图 7 给出离子和电子电场沉积率分别为 0.71±0.15 与 0.35±0.06 兆瓦，约为二比一；两物种的热能变化项却分别约为 0.27±0.10 与 0.26±0.10 兆瓦。后者还包含热储能变化和边界焓流差。表一改变控制体后，离子和电子占输入磁能的份额也发生变化，说明该比例与统计区域有关。

\paragraph{成立条件与剩余缺口}预算中的热损失具有较大误差，部分离子损失只给出上限；输出能流、边界条件和采样尺度都会影响份额。粒子模拟虽然包含真实质量比，也未覆盖任意导场、等离子体压强、三维湍流及宏观装置。本文承接 2014 年实验成果，不能把两篇文章当成相互独立的复现实验；原始数据和模拟代码未在本次复算。

\paragraph{对Hall比较的意义}常被引用的二比一是物种电场能量沉积率比，既不是离子与电子温度比，也不是两物种不可逆熵产生热比。本文证明在指定实验及控制体内可以建立有用的分物种能量预算，但没有提供普适加热闭合。应同时报告有序流动、热储能、焓流和热损失，不能只摘取一个比值作为霍尔模型的普遍预测。

\section{F679：2014年}\label{paper:F679}

\textbf{Bulk ion acceleration and particle heating during magnetic reconnection in a laboratory plasma}\citep{F679}

\textsf{本轮状态：}已有全文记录复用，不增加计数。

\textsf{核读范围：}所记录版本全部13个PDF物理页；正文、方法、结果、讨论、图表、参考文献及适用文内附录均核读。\par\textsf{图式：}直接查看PDF物理页5, 8, 9, 10, 11；Figs 5,  9–13;  Eqs (1)–(19)

\href{run:../../papers/F679.pdf}{本地原文PDF}\quad\href{https://doi.org/10.1063/1.4874331}{DOI/版本入口}

\paragraph{实际模型与问题}本文是磁重联实验装置的氦等离子体实验，并以二维全动理学粒子模拟进行对照。电子动量平衡用于解释面内电势，光谱与探针分别约束离子温度和流速；模拟比较有、无离子碰撞的速度分布。文中的“双流体区域”指电子与离子解耦的实验状态，并不意味着所用数值模型是有限矩双流体模型。

\paragraph{参数、初边值与实现}实验气压为 4.5 毫托，磁通芯间距 42 厘米，离子皮肤深度约 8 厘米，无外加导场。4200 次放电约三成通过筛选，每点平均 7 至 15 次；离子光谱的空间、时间分辨率约为 3 至 4 厘米和 5.6 微秒。粒子模拟质量比为 400，初始两物种等温，电子热速为光速的 0.183 倍；离子碰撞频率及自由程相对特征尺度与实验匹配。

\paragraph{求解得到什么}面内电场约 700 伏每米，高于约 200 伏每米的重联电场，主要离子增能由面内电势驱动。有序流能小于电势能下降量，下游出现更高离子温度。无碰撞模拟的较大温度包含多束混合；加入实验量级碰撞后，分布更接近麦克斯韦形态。电子能量沉积约 340 千瓦，经典电阻贡献估计仅约 59 千瓦。

\paragraph{成立条件与剩余缺口}实验没有完整测得电子压力张量及热通量，约 300 千瓦的热损失估计具有较大误差，因此异常电子耗散是条件性推断，不能唯一归因于所见波动。二维模拟的部分反向离子束也未被实验空间分辨率辨认。完整网格、时间步和边界细节部分依赖前序方法文献，本次未复核全部前序实现或复算原始数据。

\paragraph{对Hall比较的意义}本文把电势加速与下游分布热化联系起来，支持轨道运动、再磁化和碰撞共同作用的解释，为流体闭合提供实验约束。它没有证明只加入霍尔项就能得到正确熵产，也没有给出普适离子与电子不可逆加热比。应把物种电场功、流动能、压力功和热输运分开核算；本次依据正式版全部 13 页及关键图式作出判断。

\section{Yoo2017Heating：2017年}\label{paper:Yoo2017Heating}

\textbf{Electron heating and energy inventory during asymmetric reconnection in a laboratory plasma}\citep{Yoo2017Heating}

\textsf{本轮状态：}已有全文记录复用，不增加计数。

\textsf{核读范围：}所记录版本全部18个PDF物理页；正文、方法、结果、讨论、图表、参考文献及适用文内附录均核读。\par\textsf{图式：}直接查看PDF物理页7, 8, 10, 12, 14, 15；Figs4, 5, 8, 10, 11;  Table1;  Eqs5–10；2026-11-13纠正定位：本地PDF图10、11在第15页，式(10)在第14页，图8在第12页。

\href{run:../../papers/Yoo2017Heating.pdf}{本地原文PDF}\quad\href{https://doi.org/10.1002/2017JA024152}{DOI/版本入口}

\paragraph{实际模型与问题}本文研究两侧密度不同的磁重联实验，并使用二维和三维全动理学粒子模拟解释电子加热及物种能量收支。实验重点测量低密度侧分离面附近的电场、温度和波动；模拟比较温度矩与速度分布形状的变化。因此它既约束空间非对称条件下的能量分配，也检验分布趋于麦克斯韦形态是否必然伴随额外净增能。

\paragraph{参数、初边值与实现}实验两侧密度比约为 8，重联磁场分别约为 145 与 168 高斯，残余导场小于重联场的一成，离子皮肤深度约 7.5 厘米。粒子模拟质量比为 900、密度比为 10，初始两物种等温，高密度侧热压与磁压比为 0.5。二维盒长为 10 乘 10 个离子皮肤深度、网格 1024 平方；三维第三方向长 3 个离子皮肤深度，增加 320 层网格。

\paragraph{求解得到什么}低密度侧分离面具有增强的电子加热。图 10 中电子和离子沉积率分别为 0.35±0.07 与 0.44±0.09 兆瓦，占输入磁能约百分之二十五和三十一；两物种热能变化项均约 0.30 兆瓦。二维与三维模拟温度剖面相近，三维电子分布却更接近麦克斯韦形态，说明分布热化与净能量增长可以表现不同。

\paragraph{成立条件与剩余缺口}实验波动与温升的相关系数约为 0.29，显著性水平约为 0.086，不能确立低混杂漂移波的强因果作用。模拟的电子温度、频率比和有限盒尺度与实验仍有差异。正文温度矩定义及表格密度单位存在排印疑点，应按标准定义和正文量级核对，不能据此反推代码错误。原始数据入口已记录，但本次未重新分析。

\paragraph{对Hall比较的意义}非对称实验下两物种沉积率相近，直接限制了对称实验约二比一结果的外推。三维波动可能促进分布散射而不显著增加温度，这提示流体闭合应分别接受能量预算和分布诊断检验。本文支持针对具体条件研究物种分配，但未建立跨装置、跨采样位置的普适比例；正式版 18 页及能量收支图已完整核读。

\section{Che2014Buneman：2014年}\label{paper:Che2014Buneman}

\textbf{Two-Fluid Description of Wave-Particle Interactions in Strong Buneman Turbulence}\citep{Che2014Buneman}

\textsf{本轮状态：}已有全文记录复用，不增加计数。

\textsf{核读范围：}1-26（全部）\par\textsf{图式：}PDF7、16-19、21页视觉核验图1、6-9和完整流体方程；所有26页含图注/参考文献读完，无附录。

\href{run:../../papers/Che2014Buneman.pdf}{本地原文PDF}\quad\href{https://doi.org/10.1063/1.4882677}{DOI/版本入口}

\paragraph{实际模型与问题}3D PIC诊断与近似双流体闭合分析；未独立演化验证纯双流体重联，模拟未发展重联。

\paragraph{参数、初边值与实现}mi/me100, w0=.1di=de, Te=Ti=.04miVA\textsuperscript{2}；盒1×1×2di，512×512×1024；x/z周期、y导电；电子漂移9VA≈3vte、离子.9VA；|B|≈sqrt26 B0。

\paragraph{求解得到什么}Buneman拖曳和相混把电子漂移能转为平行热，磁能下降伴电子垂直热增长；局地流体矩预算近似刻画PIC；离子主要协助波交换，弱热化。

\paragraph{成立条件与剩余缺口}Dez(t)仍需动力学理论/PIC拟合；同前序Che2013运行，非独立验证；无实际重联、无实验比；式19为电子平行/垂直比。图7约化PIC熵增不等于精细Vlasov熵增，未熵粗粒化收敛验证。

\paragraph{对Hall比较的意义}模型身份、温度/能量/熵产区别及跨文献判断见独立加热专题；不把本篇单独作为普适分配定律。

\chapter{访问、版本及未核读线索}

Ohia2012的MIT公开记录已取得，PDF入口返回HTTP 429；OSTI旧入口未返回匹配PDF。Ohia2015出版PDF入口返回403，网页只用于检索和追引，尚未完成原文图式核验；二者不计新全文，不作为本专题科学结论的直接证据。其模型能力仍不能仅由后继论文转述代替原文核读。

Kaptanoglu2020勘误与TenBarge2019MMS正式PDF本轮访问仍未成功，沿用已读接收稿或arXiv版及已有正式网页差异记录。十三/二十矩相关会议摘要、Ohia学位论文和前序EoS原始推导保留为追引；不声称已读、已验证或已复算。

Wetherton2019使用20页OSTI接收稿，后部含重复/放大图页；Wetherton2021使用23页OSTI接收稿。两篇均无文内附录。独立仪器论文、原始MMS数据及前序方法代码未在本轮完整重新分析。